\documentclass[12pt,a4paper]{report}

\usepackage[utf8]{inputenc}
\usepackage[T1]{fontenc}
\usepackage[spanish]{babel}
\usepackage{listings}
% \usepackage{xcolor}
% \usepackage{geometry}
\usepackage{hyperref}
\usepackage{parskip}
\usepackage{titlesec}
\usepackage{enumitem}
\usepackage{booktabs}
\usepackage{tcolorbox}
\usepackage{amsmath}
\usepackage{array}

% \geometry{margin=2.5cm}

\usepackage[dvipsnames]{xcolor}
\usepackage[lmargin=2.5cm, rmargin=5cm]{geometry}
\input{../../../../../word-comments.tex}


% ─── Colores ──────────────────────────────────────────────────────────────────
\definecolor{lisp-bg}{RGB}{248,248,252}
\definecolor{lisp-kw}{RGB}{0,0,180}
\definecolor{lisp-mac}{RGB}{0,120,0}
\definecolor{lisp-str}{RGB}{160,0,0}
\definecolor{lisp-cmt}{RGB}{110,110,110}
\definecolor{box-blue}{RGB}{240,248,255}
\definecolor{box-blue-border}{RGB}{100,149,237}
\definecolor{box-green}{RGB}{240,255,240}
\definecolor{box-green-border}{RGB}{60,150,60}
\definecolor{box-yellow}{RGB}{255,253,230}
\definecolor{box-yellow-border}{RGB}{180,150,0}

% ─── Listings ─────────────────────────────────────────────────────────────────
\lstdefinelanguage{DSL}{
  keywords={def-table, libro, hoja, tabla, load, list, defparameter,
            conditional-rendering, exists, nav, ultima-fila, primera-fila,
            xl-generate, xl-run-generated},
  morekeywords=[2]{countif, counta, sum-range, str, col, trange,
                   _equals, _and, _or},
  keywordstyle=\color{lisp-kw}\bfseries,
  keywordstyle=[2]\color{lisp-mac},
  stringstyle=\color{lisp-str},
  commentstyle=\color{lisp-cmt}\itshape,
  morestring=[b]",
  morecomment=[l]{;;},
  sensitive=false,
  basicstyle=\ttfamily\small,
  backgroundcolor=\color{lisp-bg},
  frame=single,
  framerule=0.4pt,
  rulecolor=\color{gray!40},
  breaklines=true,
  numbers=left,
  numberstyle=\tiny\color{gray},
  numbersep=8pt,
  xleftmargin=20pt,
  xrightmargin=4pt,
  aboveskip=6pt,
  belowskip=6pt,
  showstringspaces=false,
  tabsize=2,
}

\lstset{language=DSL}

% ─── Cajas ────────────────────────────────────────────────────────────────────
\tcbuselibrary{skins,breakable}

\newtcolorbox{infobox}[1]{
  colback=box-blue, colframe=box-blue-border,
  fonttitle=\bfseries, title=#1,
  breakable, sharp corners, left=6pt, right=6pt, top=4pt, bottom=4pt
}

\newtcolorbox{decisionbox}[1]{
  colback=box-yellow, colframe=box-yellow-border,
  fonttitle=\bfseries, title=#1,
  breakable, sharp corners, left=6pt, right=6pt, top=4pt, bottom=4pt
}

\newtcolorbox{sintaxisbox}[1]{
  colback=box-green, colframe=box-green-border,
  fonttitle=\bfseries\ttfamily, title=Sintaxis: \texttt{#1},
  breakable, sharp corners, left=6pt, right=6pt, top=4pt, bottom=4pt
}

% ─── Títulos ──────────────────────────────────────────────────────────────────
\titleformat{\chapter}[hang]{\large\bfseries}{\thechapter.}{0.5em}{}[\vspace{2pt}\hrule]
\titleformat{\section}[hang]{\normalsize\bfseries}{\thesection}{0.5em}{}
\titleformat{\subsection}[hang]{\normalsize\itshape\bfseries}{\thesubsection}{0.5em}{}

\hypersetup{hidelinks}
\setlength{\parskip}{6pt}

% ──────────────────────────────────────────────────────────────────────────────
\begin{document}

\chapter{Tutorial: construcción de un libro de cálculo con el DSL}
\label{cap:tutorial}

\notaparaelautor{Te propongo que no pongas en el título que es para un libro de cálculo. debería ser para definir una visualización de horarios, y casi \textit{de casualidad} es que lo vamos a exportar a un libro de cálculo, pero no tiene por qué ser así.}

Este capítulo narra, paso a paso, el proceso de construir un programa completo
en el DSL a partir de la descripción de un problema de horario.\agregaesto{NUEVO PÁRRAFO.}


El ejemplo \cambio{es}{será} un \textbf{horario de defensas de tesis}:
dado un conjunto de profesores, estudiantes y locales, se necesita
\cambio{un libro de calculo}{una aplicación} que liste las defensas
programadas, \comment{compute}{calcule} automáticamente cuántas defensas acumula cada
profesor y resalte en rojo a los profesores con conflicto de horario.

La filosofía del sistema es \comment{\textbf{separar la descripción
    del problema de su generación}}{justamente por esto te propongo no
  hacer referencias al libro de cálculo.}: el programador declara la
estructura \cambio{del libro}{de la aplicación} (tablas, fórmulas, coloreado condicional) \comment{y el
generador produce el archivo \texttt{.xlsx} de forma automática.}{Y esto ponlo como una de las opciones.}

\agregaesto{Pon una estructura de las secciones.}

% ══════════════════════════════════════════════════════════════════════════════
\section{Requisitos}
\label{sec:requisitos}

\notaparaelautor{Cosa empezando con subcosa.}

\subsection{Software necesario}

Para utilizar el generador se necesitan dos programas:

\textbf{Steel Bank Common Lisp (SBCL)} \quad Intérprete de Common Lisp.
Disponible en \url{https://www.sbcl.org/}.
Versión mínima recomendada: 2.0.

\textbf{Python 3.10 o superior} \quad Para ejecutar el código generado.
Además de Python estándar, se requiere la biblioteca \textbf{openpyxl}:

\begin{lstlisting}[language=bash, numbers=none, backgroundcolor=\color{gray!10},
                   basicstyle=\ttfamily\small, keywordstyle={}]
pip install openpyxl
\end{lstlisting}

\subsection{Estructura del proyecto}

El código fuente del generador se organiza en los siguientes archivos:

\begin{lstlisting}[language={}, numbers=none, backgroundcolor=\color{gray!10},
                   basicstyle=\ttfamily\small]
generador-excel/
+-- dsl-directo.lisp         <- DSL y macros (no modificar)
+-- ast-def.lisp             <- definicion de nodos AST (no modificar)
+-- generate-code-direct.lisp <- backend de generacion (no modificar)
+-- code-generation-utils.lisp
+-- hoja_con_formulas.py     <- backend Python (no modificar)
\-- mi-horario/
    +-- data-tutorial.lisp   <- datos concretos (el usuario crea este)
    \-- tutorial-horario.lisp <- programa DSL (el usuario crea este)
\end{lstlisting}

Los archivos de la raíz son la \textbf{biblioteca del generador} y no
deben modificarse.  El usuario únicamente crea su propio programa DSL
y su archivo de datos.

\subsection{Ejecución}

\notaparaelautor{Esto déjalo para el final... en este momento del tutorial crea más ruido de lo que aclara. Eso vamos a ponerlo cuando ya tengamos todos los ficheros listos.}

Hay dos formas equivalentes de ejecutar el flujo completo.

\textbf{Opción A — script incluido} (recomendada para comenzar):

\begin{lstlisting}[language=bash, numbers=none, backgroundcolor=\color{gray!10},
                   basicstyle=\ttfamily\small, keywordstyle={}]
cd tutorial/
chmod +x generar.sh   # solo la primera vez
./generar.sh
\end{lstlisting}

\textbf{Opción B — carga directa con SBCL}:

\begin{lstlisting}[language=bash, numbers=none, backgroundcolor=\color{gray!10},
                   basicstyle=\ttfamily\small, keywordstyle={}]
cd tutorial/
sbcl --script tutorial-horario.lisp
\end{lstlisting}

Ambas opciones cargan el DSL, construyen el árbol AST, generan el programa
de construcción del libro y lo ejecutan para crear el archivo \texttt{.xlsx}.
La diferencia es que el script imprime mensajes de progreso y verifica que
el archivo fue creado correctamente.  \comment{La sección~\ref{sec:generar} explica
el flujo con más detalle.}{Perfecto... entonces vamos a dejarlo para allá, para esa sección.}

\agregaesto{TRANSICIÓN A LO QUE SIGUE.}

% ══════════════════════════════════════════════════════════════════════════════
\section{Conceptos fundamentales}
\label{sec:conceptos}

\subsection{El árbol de sintaxis abstracta (AST)}

Cuando el usuario escribe definiciones en Common Lisp usando las macros \comment{del
sistema}{¿de qué sistema?}, estas construyen automáticamente objetos en memoria que forman un
\textbf{árbol de sintaxis abstracta (AST)}.  Cada nodo del árbol representa
un concepto del problema: una tabla, una expresión de celda, un cuantificador
existencial, una regla de coloreado.

\comment{El usuario \textbf{nunca manipula el árbol directamente}: trabaja con las
macros de alto nivel descritas en este documento.}{¿por qué esto es relevante?}

\agregaesto{TRANSICIÓN A LO QUE SIGUE.}

\subsection{El DSL es independiente del generador de código}

El AST no contiene ninguna referencia a Excel, a Python \comment{ni a ninguna otra
tecnología concreta}{¿tecnologías de qué?}.  Es una representación \textbf{neutral del problema}.
\queesesto{Cualquier componente que sepa recorrer el árbol} e \queesesto{interpretar sus nodos}
puede actuar \queesesto{como generador de código} \queesesto{para el destino que desee}.

En este trabajo el generador implementado toma el AST y produce un
\textbf{script Python} que usa la biblioteca \texttt{openpyxl} para
construir \cambio{el}{un} archivo \texttt{.xlsx}.  Pero el mismo DSL podría alimentar,
sin cambio alguno en el programa del usuario, un generador diferente que
produjera, por ejemplo, una hoja LibreOffice, un informe PDF o una tabla
HTML.

\begin{infobox}{El generador implementado en este trabajo}
El backend concreto que se distribuye con el DSL genera un script Python
que usa \texttt{openpyxl} para construir archivos Excel (\texttt{.xlsx}).
Las fórmulas de celda, los datos y las reglas de formato condicional quedan
incrustados en el archivo y funcionan en cualquier versión de Excel o
LibreOffice Calc sin necesidad de tener instalado Python ni SBCL.
\end{infobox}

\subsection{Etapas de la canalización}

\notaparaelautor{¿Qué es una canalización?}

El proceso de generación tiene tres etapas \quitaesto{claramente} separadas:

\begin{enumerate}
    \item \cambio{textbfExpansion de macros.}{\textbf{Construcción del
      AST.}}  Cada \queesesto{forma del DSL} (\texttt{def-table}, \texttt{exists},
  \texttt{conditional-rendering}) \queesesto{se expande} en un nodo del AST.

    \item \textbf{Generación de código.}  \comment{El árbol}{¿qué
    árbol?} se pasa al generador, que recorre cada nodo y produce
  \comment{el programa de construcción del libro}{Sería de la
    aplicación de horarios en el formato deseado, ¿no?}.  En \comment{el
  backend incluido}{¿incluido dónde o en qué?}, ese programa es un script Python.

  \item \textbf{Ejecución.}  El programa generado se ejecuta y crea el
    archivo de salida con datos, fórmulas y reglas de formato condicional
    incrustadas.
\end{enumerate}

\agregaesto{TRANSICIÓN A LO QUE SIGUE.}

\subsection{Principio central: intención vs.\ implementación}

El DSL declara \emph{qué} debe ocurrir (``colorear los profesores con
conflicto de horario''); el generador decide \emph{cómo} implementarlo
según el destino (para el backend Excel, genera una \queesesto{\texttt{FormulaRule}
\texttt{SUMPRODUCT}} \queesesto{que persiste y recalcula en la hoja}).  El programador
del DSL no escribe ninguna fórmula ni ningún \comment{\queesesto{código del destino}}{:-o... no tiene nada que ver con la tesis, pero tiene nombre de novela brasileña :-o... Si la escribo y gano dinero por eso, te doy una parte 8-D.}.

\agregaesto{TRANSICIÓN A LO QUE SIGUE.}

% ══════════════════════════════════════════════════════════════════════════════
\section{Del problema a las tablas}
\label{sec:modelado}

El primer paso es decidir cómo representar el dominio del problema en filas
y columnas. \notaparaelautor{Perfecto... ¿por qué? Aquí sería bueno justificar por qué.}

\subsection{El problema}

\notaparaelautor{Esta no lo voy a poder escribir :-(... Arjona se me adelantó :-/.}

\notaparaelautor{Sé más específico con el nombre de la sección. Por ejemplo, para que quede claro que es \textit{el problema de horario} para el que nos intersa construir la aplicación.}

Se necesita planificar las defensas de tesis de una facultad.  Cada defensa
tiene un estudiante, un tribunal de cinco profesores (tutor, oponente,
presidente, vocal y secretario), un día, una hora y un local.  Además, se
quiere saber cuántas defensas acumula cada profesor y detectar si algún
profesor tiene dos defensas programadas a la vez.

\agregaesto{TRANSICIÓN A LO QUE SIGUE.}

\subsection{Decidir cuántas tablas hacen falta}

\begin{decisionbox}{Decisión de diseño}
La pregunta clave es: \emph{¿qué es una fila?}

Si una fila es \textbf{una defensa}, la tabla tiene tantas filas como
defensas \agregaesto{se deben planificar} y sus columnas son: estudiante, los cinco roles del tribunal,
día, hora y local.

Si una fila es \textbf{un profesor}, la tabla tiene tantas filas como
profesores \agregaesto{existen} y necesita una columna calculada que cuente en cuántas
defensas participa. \notaparaelautor{¿Y aquí no hay que agregar columnas para el nombre, el grado y eso?}

{El libro} necesita \emph{ambas} perspectivas\quitaesto{:}\agregaesto{PUNTO.} \agregaesto{Por ese motivo} se usan dos tablas en la
misma hoja {que se referencian entre sí}.
\end{decisionbox}

\notaparaelautor{Te sugiero no hablar de libro, sino de aplicación.}

\notaparaelautor{¿qué es lo que se referencian entre sí? queda ambiguo}

\notaparaelautor{De los más bonitillos los cuadraditos esos \texttt{:-D}. Me los voy a robar para los documentos de las clases \texttt{8-)}.}

\agregaesto{TRANSICIÓN A LO QUE SIGUE.}

\subsection{Tabla de defensas}

\notaparaelautor{¿Y por qué estamos hablando ahora de tablas?... ¿qué pasó con las filas?}

La \textbf{tabla de defensas} (\texttt{defensas-table}) \comment{tiene una fila por defensa}{Abusador... pero si lo que me vas a poner ahora son las columnas, ¿por qué me hablas de las filas? :-/}:

\begin{center}
\begin{tabular}{ll}
\toprule
\textbf{Columna}     & \textbf{Descripción} \\
\midrule
\texttt{estudiante}  & Nombre del estudiante que defiende \\
\texttt{tutor}       & Profesor tutor \\
\texttt{oponente}    & Profesor oponente \\
\texttt{presidente}  & Presidente del tribunal \\
\texttt{vocal}       & Vocal del tribunal \\
\texttt{secretario}  & Secretario del tribunal \\
\texttt{dia}         & Día de la defensa \\
\texttt{hora}        & Hora de inicio \\
\texttt{local}       & Nombre del local \\
\bottomrule
\end{tabular}
\end{center}

\begin{infobox}{Por qué los cinco roles van juntos y en ese orden}
  Los roles \texttt{tutor}, \texttt{oponente}, \texttt{presidente},
  \texttt{vocal} y \texttt{secretario} son columnas consecutivas.
  Esta disposición permite referenciarlos como {un rango
    continuo}\footnote{¿Qué es un rango continuo?} {\texttt{(trange defensas-table tutor
      secretario)}} {al calcular cuántas veces aparece un
    profesor en cualquier rol}, sin necesidad de {sumar cinco
  búsquedas separadas.}\footnote{¿Por qué hay que sumar 5?}
\end{infobox}

\agregaesto{TRANSICIÓN A LO QUE SIGUE.}

\notaparaelautor{¿Y ya creamos la tabla de las defensas? ¿por qué me hablas de ella si no la vamos a crear ahora?}

\subsection{Tabla de profesores}

La \textbf{tabla de profesores} (\texttt{profesores-table}) \comment{tiene una
fila por profesor:}{Y dale con las filas... para después hablar de columnas.}

\begin{center}
\begin{tabular}{ll}
\toprule
\textbf{Columna}   & \textbf{Descripción} \\
\midrule
\texttt{nombre}    & Nombre del profesor \\
\texttt{grado}     & Grado académico \\
\texttt{defensas}  & Número de defensas en que participa (calculada) \\
\bottomrule
\end{tabular}
\end{center}

\notaparaelautor{¿Y ya? ¿se acabó esta sección? :-o... ¡qué falta de personalidad tiene esta sección! :-/. Sácale más partido.}

\agregaesto{TRANSICIÓN A LO QUE SIGUE.}

% ══════════════════════════════════════════════════════════════════════════════
\section{Definir la tabla de defensas}
\label{sec:defensas-table}

\notaparaelautor{Brutal empezar la sección con el código...}

\begin{sintaxisbox}{def-table}
\begin{lstlisting}
(def-table nombre-tabla
  ((col1 "Cabecera1") (col2 "Cabecera2") ...)
  :cell-height N
  :cell-width  N)
\end{lstlisting}
\end{sintaxisbox}

Cada par \comment{\texttt{(columna "Cabecera")}}{revisa que cabecera
  aparece con un chirimbolo} declara \queesesto{el nombre interno}
(símbolo DSL) y\queesesto{ el texto del encabezado visual de la
  columna}.  Las opciones \texttt{:cell-height} y \texttt{:cell-width}
controlan cuántas filas y columnas físicas \comment{de Excel}{¿no
  habíamos quedado en que esta primera parte era agnóstica al
  \textit{backend}?} ocupa cada \queesesto{fila o columna lógica} (el
valor 1 es el caso habitual).

\agregaesto{TRANSICIÓN A LA TABLA QUE SIGUE.}

\begin{lstlisting}[caption={Definicion de defensas-table.}]
(def-table defensas-table
  ((estudiante "") (tutor "") (oponente "") (presidente "")
   (vocal "") (secretario "") (dia "") (hora "") (local ""))
  :cell-height 1
  :cell-width 1)
\end{lstlisting}

\comment{El valor \texttt{""} en cada par indica que la columna se rellena con datos
externos}{Aquí te propondría usar otro símbolo para eso, que en definitiva los símbolos nos salen gratis.}; \queesesto{no es un encabezado visible} {sino el valor por defecto de la celda
antes de que el generador coloque las fórmulas.}

\notaparaelautor{Me pasa lo mismo que en la versión anterior... deja esta tabla para lo último... después que tengas las demás cosas definidas.}

\notaparaelautor{Y otra cosa... ¿eso no tengo que agregarlo a algún fichero?}

\agregaesto{TRANSICIÓN A LO QUE SIGUE.}

% ══════════════════════════════════════════════════════════════════════════════
\section{Definir la tabla de profesores}
\label{sec:profesores-table}

La tabla de profesores introduce tres características nuevas: columnas
calculadas, referencias cruzadas y formato condicional. \agregaesto{Justifica aquí por qué esa tabla necesita esas funcionalidades.}

\agregaesto{TRANSICIÓN A LO QUE SIGUE.}

\subsection{Columna calculada con \texttt{:computed}}

La opción \texttt{:computed} asocia columnas a \queesesto{expresiones
  DSL} \comment{que el generador traduce a fórmulas Excel}{Esto es
  demasiado específico}.  El valor de esa celda en los datos de
entrada se ignora: la celda mostrará siempre el resultado de la
fórmula.

\agregaesto{Pon un ejemplo con palabras antes de poner el código.}

\agregaesto{TRANSICIÓN AL CÓDIGO.}

\begin{lstlisting}[caption={Columna defensas calculada con COUNTIF.}]
(def-table profesores-table
  ((nombre "") (grado "") (defensas ""))
  :cell-height 1
  :cell-width  1
  :computed
    ((defensas (countif (trange defensas-table tutor secretario)
                        (col nombre)))))
\end{lstlisting}

La expresión \texttt{(countif rango criterio)} \comment{genera la fórmula Excel
\texttt{COUNTIF(rango, criterio)}}{Demasiado específico de Excel, esta parte del tutorial debería ser genérica, para que se pueda implementar en cualquier \textit{backend}}.

\agregaesto{TRANSICIÓN A LO QUE SIGUE.}

\begin{infobox}{Cómo funciona \texttt{trange}}
  \texttt{(trange defensas-table tutor secretario)} genera un
  rango\footnote{\notaparaelautor{¿qué es un rango?}} que abarca desde la columna
  \texttt{tutor} hasta la columna \texttt{secretario}\footnote{\notaparaelautor{¿Y si quisiera buscar en varias columnas que no sean consecutivas?}} de
  \texttt{defensas-table}, para todas las filas de datos.  En Excel\footnote{\notaparaelautor{Eso es demasiado específico para este momento del tutorial. Deberías explicar el funcionamiento de manera genérica}} se
  traduce a algo como \texttt{\$B\$4:\$F\$6}.

La expresión completa\footnote{\notaparaelautor{¿qué expresión completa? Aquí solo estás hablando del trange}} cuenta cuántas veces aparece el nombre del profesor
actual (tomado de su propia columna \texttt{nombre} con \texttt{(col nombre)})
en cualquiera de los cinco roles de cualquier defensa.

Se usa \texttt{trange} en lugar de \texttt{range}\footnote{\notaparaelautor{¿quéjesto? :-o}} porque ambas tablas
tienen una columna llamada \texttt{nombre}: \texttt{trange} especifica
explícitamente la tabla de origen y evita la ambigüedad. \notaparaelautor{¿Y entonces cuál es la diferencia conceptual entre trange y range?}
\end{infobox}

\agregaesto{TRANSICIÓN A LO QUE SIGUE.}

\subsection{Formato condicional con \texttt{:render}}

La opción \texttt{:render} acepta una forma \queesesto{\texttt{conditional-rendering}}
que especifica bajo qué condición se colorean determinadas columnas.
La condición se expresa con el cuantificador \texttt{exists}.

\agregaesto{Primero pon un ejemplo con palabras... antes del código.}

\begin{lstlisting}[caption={Formato condicional sobre profesores-table.}]
(def-table profesores-table
  ((nombre "") (grado "") (defensas ""))
  :cell-height 1
  :cell-width  1
  :computed
    ((defensas (countif (trange defensas-table tutor secretario)
                        (col nombre))))
  :render
    (conditional-rendering
      :condition
        (exists (r :rows-of defensas-table)
          :self-in  (tutor oponente presidente vocal secretario)
          :matching (dia hora))
      :target-columns (nombre)))
\end{lstlisting}

La variable \texttt{r} es \queesesto{un alias documental} de la \queesesto{fila que se comprueba},
análogo a un alias de tabla en SQL.

\agregaesto{Ejemplo aquí... que no entiendo mucho...}

\begin{sintaxisbox}{exists :rows-of}
\begin{lstlisting}
(exists (var :rows-of tabla-id)
  :self-in  (col1 col2 ...)
  :matching (clave1 clave2 ...))
\end{lstlisting}
\end{sintaxisbox}

\begin{center}
\renewcommand{\arraystretch}{1.3}
\begin{tabular}{>{\ttfamily}p{3.5cm}p{8.5cm}}
\toprule
\textbf{Parámetro} & \textbf{Significado} \\
\midrule
\texttt{:rows-of} & Tabla en la que se busca \\
\texttt{:self-in} & Lista de columnas de esa tabla donde debe aparecer
                    el valor de la celda actual \\
\texttt{:matching} & Columnas cuyo slot debe tener más de una defensa
                     programada (conflicto) \\
\bottomrule
\end{tabular}
\end{center}

La semántica de la condición del ejemplo es: ``existe alguna fila de
\texttt{defensas-table} en la que el nombre de este profesor aparece en
alguno de los cinco roles, y cuyo par (\texttt{dia}, \texttt{hora}) tiene
más de una defensa programada''.

\begin{infobox}{El coloreado persiste en Excel}
\notaparaelautor{Esto es demasiado específico de excel... en este punto del tutorial esos detalles no inmportan... de hecho, si la persona no sabe excel, no sabe qué está mirando aquí. Y hay muuuuchas personas que no saben excel :-(.}
  
El generador no colorea celdas directamente: incrusta una
\texttt{FormulaRule} en el archivo \texttt{.xlsx}.
Esta regla se reevalúa automáticamente cada vez que el usuario edita los
datos en Excel, sin necesidad de volver a ejecutar el DSL.

La fórmula generada para este ejemplo es del tipo:
\begin{lstlisting}[numbers=none, backgroundcolor=\color{gray!10},
                   basicstyle=\ttfamily\footnotesize, language={}]
=SUMPRODUCT(
    (($B$4:$B$6=$K4)+...+($F$4:$F$6=$K4)) *
    (COUNTIFS($G$4:$G$6,$G$4:$G$6,$H$4:$H$6,$H$4:$H$6)>1)
) > 0
\end{lstlisting}
\end{infobox}

\agregaesto{TRANSICIÓN A LO QUE SIGUE.}

% ══════════════════════════════════════════════════════════════════════════════
\section{Preparar los datos}
\label{sec:datos}

Los datos se guardan en un archivo separado como parámetros globales de Lisp.
Cada parámetro es una lista de listas: una sublista por fila.

\subsection{Formato de los datos}

Cada bloque en la hoja de cálculo tiene \queesesto{\textbf{tres filas de prefijo}}
seguidas de las filas de datos. \comment{ Las filas de prefijo se escriben tal cual
(no generan fórmulas)}{¿y por qué el nombre prefijo? ¿dónde se escriben?}: la primera puede llevar \comment{un identificador opcional}{¿para qué?},
la segunda está vacía y la tercera contiene los encabezados de columna.

\notaparaelautor{quizás sea bueno incluir en el tutorial ejemplos de tablas (de LaTeX de toda la vida) y a partir de ellas ir mostrando el lenguaje. Por ejemplo: vamos a escribir esta tablita... que no hace nada. se hace así. Ahora vamos a incorporarle un formato condicional... para que pase tal cosa... se hace así... y ahora vamos a ponerle un fórmula sencillita, para calcular la suma de los dos... mira qué lindo se hace. Quizás eso debería ser la primera parte del tutorial... poder definir tablitas ahí satas, sin grandes pretensiones... y después, cuando el usuario haya generado un librito bobo, entrarle entonces a las cosas del horario... porque creo que las dos cosas a la vez está complicado... :-/... la dos cosas las cosas de horario y la sintaxis del lenguaje.}

\notaparaelautor{Y esto mismo de los encabezados, los datos y esas cositas, las presentamos con las tablas bobas...}

\begin{lstlisting}[caption={Datos de defensas-table.}]
(defparameter *DATOS-DEFENSAS*
  '(("" "" "" "" "" "" "" "" "")         ;; fila 1: prefijo
    ("" "" "" "" "" "" "" "" "")         ;; fila 2: separador
    ("Estudiante" "Tutor" "Oponente"     ;; fila 3: encabezados
     "Presidente" "Vocal" "Secretario" "Dia" "Hora" "Local")
    ("Juan Diaz"  "Piad"   "Garcia" "Lopez"  "Torres" "Ruiz"
                                         "Lunes"  "10:00" "Aula 1")
    ("Maria Vega" "Garcia" "Lopez"  "Piad"   "Soto"   "Torres"
                                         "Lunes"  "14:00" "Aula 2")
    ("Luis Mora"  "Lopez"  "Torres" "Garcia" "Ruiz"   "Soto"
                                         "Martes" "10:00" "Aula 1")))
\end{lstlisting}

Las celdas con fórmulas calculadas \comment{se dejan vacías \texttt{""}}{Quizás aquí tenga más sentido todavía usar otro símbolo que no sea la cadena vacía... fíjate, a lo mejor es una variable de lisp que se \textup{formula} y que lo que tenga es la cadena vacía.} en el
archivo de datos: el generador las sobreescribe con la fórmula
correspondiente.



\begin{lstlisting}[caption={Datos de profesores-table.}]
(defparameter *DATOS-PROFESORES*
  '(("" "" "")
    ("" "" "")
    ("Nombre" "Grado" "Defensas")
    ("Piad"   "Dr"    "")     ;; "Defensas" se calcula con COUNTIF
    ("Garcia" "Dr"    "")
    ("Lopez"  "Msc"   "")
    ("Torres" "Msc"   "")
    ("Ruiz"   "Lic"   "")
    ("Soto"   "Msc"   "")))
\end{lstlisting}

\agregaesto{TRANSICIÓN A LO QUE SIGUE.}

% ══════════════════════════════════════════════════════════════════════════════
\section{Construir el libro}
\label{sec:libro}

Con las tablas declaradas y los datos preparados, se ensambla el \queesesto{workbook}
con tres macros: \texttt{libro}, \texttt{hoja} y \texttt{tabla}.

\begin{sintaxisbox}{libro / hoja / tabla}
\begin{lstlisting}
(libro nombre-workbook
  :filename "nombre.xlsx"
  :hojas (list
    (hoja "Nombre hoja"
      (tabla plantilla1 :data *DATOS-1*)
      (tabla plantilla2 :data *DATOS-2*)
      :fixed-expressions
        (((nav ancla dir N) (expresion))
         ...))))
\end{lstlisting}
\end{sintaxisbox}

La macro \texttt{tabla} instancia \queesesto{una plantilla} (\texttt{def-table}) con
datos concretos.  \comment{Cuando una \texttt{hoja} recibe más de una \texttt{tabla},
el generador las coloca lado a lado dentro de la misma \queesesto{\emph{región
horizontal}}, separadas por una columna en blanco.}{Esto es un detalle del backend de excel... no es relevante todavía.}

La clave \texttt{:fixed-expressions} coloca expresiones o etiquetas fuera
de las tablas.  La posición se especifica con \texttt{nav}, \comment{que calcula la
celda}{esto es solo relevante para excel... en emacs no existe esa posiblidad porque no todo son celdas así que no estoy seguro de que debería estar en el lenguaje... } a partir de la última fila de una tabla más una serie de pasos
direccionales.

\begin{lstlisting}[caption={Workbook completo con totales de pie de tabla.}]
(libro horario-tesis
  :filename "Horario_Tesis.xlsx"
  :hojas (list
    (hoja "Defensas"
      (tabla defensas-table   :data *DATOS-DEFENSAS*)
      (tabla profesores-table :data *DATOS-PROFESORES*)
      :fixed-expressions
        (((nav (ultima-fila defensas-table estudiante) abajo 1)
          (str "Total defensas:"))
         ((nav (ultima-fila defensas-table estudiante) abajo 1 derecha 1)
          (counta (trange defensas-table estudiante)))
         ((nav (ultima-fila profesores-table nombre) abajo 1)
          (str "Total profesores:"))
         ((nav (ultima-fila profesores-table nombre) abajo 1 derecha 2)
          (sum-range (trange profesores-table defensas)))))))
\end{lstlisting}

\notaparaelautor{Aquí horario sigue teniendo un filename... eso me tiene muy nervioso...}

\agregaesto{TRANSICIÓN A LO QUE SIGUE.}

\begin{sintaxisbox}{nav / ultima-fila}
\begin{lstlisting}
;; Posicion relativa a la ultima fila de una columna:
(nav (ultima-fila tabla-id columna) abajo  N)
(nav (ultima-fila tabla-id columna) abajo  N derecha M)
(nav (primera-fila tabla-id columna) arriba N)
\end{lstlisting}
\end{sintaxisbox}

\notaparaelautor{Recuerda que en otros contextos esto no tiene por qué existir, tanto lo que está en el bloque anterior como en el que sigue.}

\begin{infobox}{Por qué usar \texttt{nav} en lugar de coordenadas fijas}
Las coordenadas absolutas de la celda ``una fila debajo de la última
defensa'' dependen de cuántas defensas haya en los datos.
\texttt{nav} las calcula automáticamente en tiempo de generación, de modo
que si se añaden más filas de datos los totales se desplazan solos sin
modificar el programa DSL\footnote{\notaparaelautor{¿qué es \textit{el programa DSL}?}}.
\end{infobox}

\agregaesto{TRANSICIÓN A LO QUE SIGUE.}
% ══════════════════════════════════════════════════════════════════════════════
\section{Generar y ejecutar}
\label{sec:generar}

Las dos últimas líneas del programa desencadenan la generación del archivo:

\begin{lstlisting}[caption={Compilacion y ejecucion desde el programa DSL.}]
(xl-generate horario-tesis "horario-tesis.py")
(xl-run-generated "horario-tesis.py")
\end{lstlisting}

\texttt{xl-generate} recorre el AST del workbook y produce el programa de
construcción del libro.
\texttt{xl-run-generated} ejecuta ese programa con \texttt{python3};
el resultado es el archivo \texttt{Horario\_Tesis.xlsx} en disco.

El archivo \texttt{.xlsx} contiene tanto las fórmulas de celda como las
\texttt{FormulaRule} de formato condicional; estas últimas permanecen
activas cuando el usuario abre y edita el libro directamente en Excel.

\subsection{Ejecución desde la terminal con el script \texttt{generar.sh}}

Como alternativa a invocar SBCL directamente, el directorio incluye un
script Bash \texttt{generar.sh} que automatiza el mismo flujo:

\begin{lstlisting}[language=bash, numbers=none, backgroundcolor=\color{gray!10},
                   basicstyle=\ttfamily\small, keywordstyle={}]
# Primera vez: dar permisos de ejecucion
chmod +x generar.sh

# Ejecutar el flujo completo
./generar.sh
\end{lstlisting}

El script realiza exactamente los mismos tres pasos que ocurren al cargar
el programa DSL a mano:

\begin{enumerate}
  \item Sitúa el directorio de trabajo en la carpeta del tutorial (para que
        los \texttt{load} relativos del programa DSL funcionen
        independientemente del directorio desde el que se llame al script).
  \item Invoca \texttt{sbcl -{}-script tutorial-horario.lisp}, que expande
        las macros, construye el AST, genera \texttt{horario-tesis.py} y lo
        ejecuta con \texttt{python3}.
  \item Verifica que \texttt{Horario\_Tesis.xlsx} fue creado e imprime su
        tamaño.
\end{enumerate}

\begin{infobox}{¿Cuándo usar el script y cuándo cargar directamente?}
El script es conveniente para ejecutar el tutorial de un solo paso desde
cualquier directorio.  Cargar el archivo directamente desde el REPL de
SBCL (\texttt{(load "tutorial-horario.lisp")}) es preferible durante el
desarrollo, porque permite inspeccionar el estado del AST en memoria entre
las distintas formas del DSL.
\end{infobox}

% ══════════════════════════════════════════════════════════════════════════════
\section{El programa completo}
\label{sec:completo}

A continuación se presenta el contenido íntegro de los dos archivos que
conforman el programa.

\subsection{\texttt{data-tutorial.lisp}}

\begin{lstlisting}[caption={data-tutorial.lisp --- datos de entrada.}]
(defparameter *DATOS-DEFENSAS*
  '(("" "" "" "" "" "" "" "" "")
    ("" "" "" "" "" "" "" "" "")
    ("Estudiante" "Tutor" "Oponente" "Presidente" "Vocal"
     "Secretario" "Dia" "Hora" "Local")
    ("Juan Diaz"  "Piad"   "Garcia" "Lopez"  "Torres" "Ruiz"
                                   "Lunes"  "10:00" "Aula 1")
    ("Maria Vega" "Garcia" "Lopez"  "Piad"   "Soto"   "Torres"
                                   "Lunes"  "14:00" "Aula 2")
    ("Luis Mora"  "Lopez"  "Torres" "Garcia" "Ruiz"   "Soto"
                                   "Martes" "10:00" "Aula 1")))

(defparameter *DATOS-PROFESORES*
  '(("" "" "")
    ("" "" "")
    ("Nombre" "Grado" "Defensas")
    ("Piad"   "Dr"    "")
    ("Garcia" "Dr"    "")
    ("Lopez"  "Msc"   "")
    ("Torres" "Msc"   "")
    ("Ruiz"   "Lic"   "")
    ("Soto"   "Msc"   "")))
\end{lstlisting}

\subsection{\texttt{tutorial-horario.lisp}}

\begin{lstlisting}[caption={tutorial-horario.lisp --- programa DSL completo.}]
(load "dsl-directo.lisp")
(load "data-tutorial.lisp")

;; TABLA 1: defensas
;; Nueve columnas: cinco roles consecutivos (tutor...secretario),
;; dia, hora, local.  Sin coloreado propio.
(def-table defensas-table
  ((estudiante "") (tutor "") (oponente "") (presidente "")
   (vocal "") (secretario "") (dia "") (hora "") (local ""))
  :cell-height 1
  :cell-width 1)

;; TABLA 2: profesores
;; Columna "defensas" calculada con COUNTIF.
;; Coloreado condicional: nombre en rojo si el profesor tiene conflicto.
(def-table profesores-table
  ((nombre "") (grado "") (defensas ""))
  :cell-height 1
  :cell-width  1
  :computed
    ((defensas (countif (trange defensas-table tutor secretario)
                        (col nombre))))
  :render
    (conditional-rendering
      :condition
        (exists (r :rows-of defensas-table)
          :self-in  (tutor oponente presidente vocal secretario)
          :matching (dia hora))
      :target-columns (nombre)))

;; WORKBOOK
(libro horario-tesis
  :filename "Horario_Tesis.xlsx"
  :hojas (list
    (hoja "Defensas"
      (tabla defensas-table   :data *DATOS-DEFENSAS*)
      (tabla profesores-table :data *DATOS-PROFESORES*)
      :fixed-expressions
        (((nav (ultima-fila defensas-table estudiante) abajo 1)
          (str "Total defensas:"))
         ((nav (ultima-fila defensas-table estudiante) abajo 1 derecha 1)
          (counta (trange defensas-table estudiante)))
         ((nav (ultima-fila profesores-table nombre) abajo 1)
          (str "Total profesores:"))
         ((nav (ultima-fila profesores-table nombre) abajo 1 derecha 2)
          (sum-range (trange profesores-table defensas)))))))

(xl-generate horario-tesis "horario-tesis.py")
(xl-run-generated "horario-tesis.py")
\end{lstlisting}

% ══════════════════════════════════════════════════════════════════════════════
\section{Resultado esperado}
\label{sec:resultado}

El programa genera \texttt{Horario\_Tesis.xlsx} con la siguiente estructura:

\begin{center}
\renewcommand{\arraystretch}{1.3}
\begin{tabular}{ll}
\toprule
\textbf{Rango} & \textbf{Contenido} \\
\midrule
Columnas A--I, filas 4--6 & \texttt{defensas-table}: tres defensas \\
Columna J                 & Separador en blanco \\
Columnas K--M, filas 4--9 & \texttt{profesores-table}: seis profesores \\
Columna M, filas 4--9     & Fórmulas \texttt{COUNTIF} generadas \\
Fila 7 (cols.\ A--B)      & Totales de defensas (\texttt{COUNTA}) \\
Fila 11 (cols.\ K, M)     & Totales de profesores (\texttt{SUM}) \\
Rango K4:K9               & \texttt{FormulaRule}: nombres en rojo
                            si hay conflicto \\
\bottomrule
\end{tabular}
\end{center}

Con los datos de ejemplo, las defensas de ``Juan Díaz'' y ``Luis Mora''
comparten el slot \texttt{Lunes 10:00} con roles solapados, por lo que
los profesores Piad, García, López, Torres, Ruiz y Soto aparecen
resaltados en rojo.

% ══════════════════════════════════════════════════════════════════════════════
\section{Traza del diseño}
\label{sec:traza}

El cuadro siguiente resume las decisiones tomadas durante el modelado,
relacionando cada requisito del problema con su representación en el DSL.

\begin{center}
\renewcommand{\arraystretch}{1.3}
\begin{tabular}{p{4.2cm}p{3.2cm}p{5.2cm}}
\toprule
\textbf{Requisito} & \textbf{Representación} & \textbf{Decisión de diseño} \\
\midrule
Una fila por defensa con tribunal completo
  & \texttt{defensas-table}
  & Granularidad natural; agrupa todas las columnas de tribunal \\
Los cinco roles en columnas consecutivas
  & orden en \texttt{def-table}
  & Permite \texttt{trange tutor secretario} para contar roles \\
Contar defensas por profesor
  & columna \texttt{computed}
  & \texttt{countif} sobre el rango de roles de \texttt{defensas-table} \\
Detectar profesores con conflicto
  & \texttt{:render} con \texttt{exists :rows-of}
  & El backend genera una \texttt{FormulaRule} que persiste en Excel \\
Totales dinámicos al pie de tabla
  & \texttt{:fixed-expressions} con \texttt{nav}
  & Las coordenadas se recalculan si cambia el número de filas \\
\bottomrule
\end{tabular}
\end{center}

% ══════════════════════════════════════════════════════════════════════════════
\section{Referencia rápida de formas}
\label{sec:referencia}

\begin{center}
\renewcommand{\arraystretch}{1.4}
\begin{tabular}{>{\ttfamily}p{5.6cm}p{7.4cm}}
\toprule
\textbf{Forma} & \textbf{Descripción} \\
\midrule
(def-table nom cols opts...)     & Define una plantilla de tabla \\
:computed ((col expr))           & Columna calculada con fórmula Excel \\
:render (conditional-rendering)  & Regla de coloreado condicional \\
(conditional-rendering :condition :target-columns)
                                 & Asocia condición con columnas a colorear \\
(exists (r :rows-of tabla) ...)  & Existencia en otra tabla (cross-table) \\
(exists (r :all-rows) ...)       & Existencia en la misma tabla \\
(col nombre)                     & Valor de la columna en la fila actual \\
(trange tabla desde hasta)       & Rango en otra tabla \\
(countif rango criterio)         & Fórmula COUNTIF \\
(counta rango)                   & Fórmula COUNTA \\
(sum-range rango)                & Fórmula SUM \\
(str "texto")                    & Literal de texto en celda \\
(nav (ultima-fila t c) dir N)    & Posición relativa a última fila \\
(nav (primera-fila t c) dir N)   & Posición relativa a primera fila \\
(libro nom :filename :hojas)     & Crea el workbook \\
(hoja nom tablas :fixed-expressions)
                                 & Define una hoja del libro \\
(tabla plantilla :data datos)    & Instancia una plantilla con datos \\
(xl-generate wb archivo)         & Genera el programa de construcción \\
(xl-run-generated archivo)       & Ejecuta el programa y crea el .xlsx \\
\bottomrule
\end{tabular}
\end{center}

% ══════════════════════════════════════════════════════════════════════════════
\section{Formas del DSL no utilizadas en este tutorial}
\label{sec:macros-extra}

El tutorial cubre un subconjunto del DSL suficiente para construir el horario
de defensas. Las formas que se describen a continuación están disponibles y
resultan útiles en escenarios más complejos.

% ─────────────────────────────────────────────────────────────────
\subsection*{Condicionales en celdas}

Permiten que el valor de una columna dependa de una condición evaluada fila
a fila.

\begin{center}
\begin{tabular}{lp{8cm}}
\toprule
\textbf{Forma} & \textbf{Para qué sirve} \\
\midrule
\texttt{(\_if cond entonces sino)} & IF de tres ramas.  Muestra un valor sólo
  cuando se cumple la condición; en caso contrario muestra otro. \\
\texttt{(non-empty expr)} & Verdad si la expresión no está vacía (\texttt{<> ""}).
  Útil como condición de \texttt{\_if}. \\
\texttt{(show-nothing)} & Produce la cadena vacía \texttt{""}.  Se usa como
  rama \emph{sino} para dejar la celda en blanco. \\
\texttt{(equals a b)} / \texttt{(different a b)} & Igualdad y desigualdad
  entre dos expresiones. \\
\texttt{(\_and a b)} / \texttt{(\_or a b)} & Conjunción y disyunción lógica. \\
\bottomrule
\end{tabular}
\end{center}

Ejemplo: mostrar el nombre de una sala sólo si la fila tiene un evento
asignado, y dejarla en blanco en caso contrario.

\begin{lstlisting}
(col "Sala" :computed
  (_if (non-empty (col "Evento"))
       (col "Sala")
       (show-nothing)))
\end{lstlisting}

% ─────────────────────────────────────────────────────────────────
\subsection*{Navegación entre filas}

Permiten que una columna acceda al valor de la misma celda en la fila
anterior o siguiente, lo que es indispensable para calcular diferencias
temporales entre registros consecutivos.

\begin{center}
\begin{tabular}{lp{8cm}}
\toprule
\textbf{Forma} & \textbf{Para qué sirve} \\
\midrule
\texttt{(previous-row expr)} & Desplaza la referencia una fila hacia arriba. \\
\texttt{(previous-of base)} & Valor de \texttt{base} en la fila inmediatamente
  anterior. \\
\texttt{(next-of base)} & Valor de \texttt{base} en la fila siguiente. \\
\texttt{(it-is-the-first-row)} & Verdad en la primera fila de datos.  Se
  combina con \texttt{\_if} para el caso borde donde no existe fila anterior. \\
\texttt{primera-fila} & Ancla de posición a la primera fila de la tabla;
  complemento de \texttt{ultima-fila}. \\
\bottomrule
\end{tabular}
\end{center}

Ejemplo: calcular la duración de cada franja horaria como la diferencia
entre su hora de fin y la hora de fin de la fila anterior.

\begin{lstlisting}
(col "Duracion" :computed
  (_if (it-is-the-first-row)
       (str "")
       (subtract (col "Hora fin")
                 (previous-of (col "Hora fin")))))
\end{lstlisting}

% ─────────────────────────────────────────────────────────────────
\subsection*{Aritmética y texto}

\begin{center}
\begin{tabular}{lp{8cm}}
\toprule
\textbf{Forma} & \textbf{Para qué sirve} \\
\midrule
\texttt{(add a b)} / \texttt{(subtract a b)} & Suma y resta sobre valores
  numéricos o de celda. \\
\texttt{(multiply a b)} / \texttt{(divide a b)} & Multiplicación y división. \\
\texttt{(time-add hora minutos)} & Suma una cantidad de minutos a una hora
  Excel (\texttt{hora + minutos/1440}). \\
\texttt{(concat a b)} & Concatenación de texto: combina el contenido de dos
  expresiones en una sola cadena. \\
\bottomrule
\end{tabular}
\end{center}

Ejemplo: calcular la hora de fin a partir de la hora de inicio y una
duración fija en minutos.

\begin{lstlisting}
(col "Hora fin" :computed
  (time-add (col "Hora inicio") (col "Duracion min")))
\end{lstlisting}

% ─────────────────────────────────────────────────────────────────
\subsection*{Referencias avanzadas entre tablas}

\begin{center}
\begin{tabular}{lp{8cm}}
\toprule
\textbf{Forma} & \textbf{Para qué sirve} \\
\midrule
\texttt{(tcol tabla nombre)} & Como \texttt{col}, pero especifica
  explícitamente de qué tabla proviene la columna.  Necesario cuando dos
  tablas de la misma hoja comparten un nombre de columna y hay que
  desambiguar. \\
\texttt{(range desde hasta)} & Rango de columnas de la tabla \emph{actual},
  sin necesidad de nombrarla.  Equivalente de \texttt{trange} para la tabla
  en curso. \\
\texttt{(lookup clave campo)} & Búsqueda estilo VLOOKUP: dado el valor de
  \texttt{clave} en la fila actual, devuelve el valor de \texttt{campo} en
  la tabla de referencia. \\
\bottomrule
\end{tabular}
\end{center}

% ─────────────────────────────────────────────────────────────────
\subsection*{Plantillas paramétricas}

Cuando una misma estructura de tabla debe instanciarse varias veces con un
detalle variable ---por ejemplo, la misma tabla de franjas para cada día de
la semana--- se declaran \emph{inst-params} en \texttt{def-table}:

\begin{lstlisting}
(def-table franja-dia
  ((turno "") (salon "") (grupo ""))
  :inst-params (dia)
  :computed
    ((salon (tcol sala-disponible salon))))
\end{lstlisting}

Dentro de la definición, \texttt{(param nombre)} accede al valor que se
pase al instanciar.  Al construir la hoja se indica el parámetro como
argumento con nombre:

\begin{lstlisting}
(hoja "Semana"
  (tabla franja-dia :dia 'lunes  :data *datos-lunes*)
  (tabla franja-dia :dia 'martes :data *datos-martes*))
\end{lstlisting}

Con este mecanismo el DSL genera cinco hojas con estructura idéntica y
fórmulas parametrizadas con el nombre del día, a partir de una sola
declaración \texttt{def-table}.

% ─────────────────────────────────────────────────────────────────
\subsection*{Detección de solapamiento dentro de la misma tabla}

El tutorial usa \texttt{(exists (r :rows-of tabla) \ldots)} para buscar
conflictos en una tabla \emph{distinta}.  La variante \texttt{:all-rows}
detecta solapamientos dentro de la \emph{propia} tabla:

\begin{lstlisting}
(conditional-rendering
  (exists (r :all-rows)
    :matching (:col-a "Inicio" :col-b "Fin")
    :any-overlap t)
  :format '(:bg "#FF9999"))
\end{lstlisting}

Útil, por ejemplo, para marcar franjas que coinciden en la misma tabla de
turnos de un único recurso.

% ─────────────────────────────────────────────────────────────────
\subsection*{Celdas de varias filas o columnas}

Las opciones \texttt{:cell-height~N} y \texttt{:cell-width~N} de
\texttt{def-table} hacen que cada fila lógica ocupe \texttt{N} filas
(o columnas) físicas en el Excel generado.  Es adecuado para horarios
donde cada franja horaria debe abarcar varias celdas de la cuadrícula,
o para encabezados que fusionan columnas.

\begin{lstlisting}
(def-table horario-visual
  ((hora "") (actividad "") (sala ""))
  :cell-height 2)   ; cada franja ocupa 2 filas fisicas
\end{lstlisting}

% ─────────────────────────────────────────────────────────────────
\subsection*{Hojas con tablas apiladas verticalmente}

La macro \texttt{hoja-v} es la versión vertical de \texttt{hoja}: apila
las tablas una encima de otra en lugar de colocarlas en columnas.
Se usa cuando las tablas tienen muchas columnas y el libro se orienta para
impresión en apaisado, o cuando las tablas representan secciones
consecutivas de un mismo listado.

\begin{lstlisting}
(hoja-v "Resumen"
  (tabla tabla-manana :data *datos-manana*)
  (tabla tabla-tarde  :data *datos-tarde*))
\end{lstlisting}

% ─────────────────────────────────────────────────────────────────
\subsection*{Referencias cruzadas entre hojas}

Para libros con varias hojas que necesitan compartir datos, el DSL ofrece
dos formas:

\begin{center}
\begin{tabular}{lp{8cm}}
\toprule
\textbf{Forma} & \textbf{Para qué sirve} \\
\midrule
\texttt{(sheet-ref hoja plantilla-celda)} & Genera una referencia a una
  celda de otra hoja, expresada como plantilla de texto en la que
  \texttt{\$\{row-num\}} se sustituye por el número de fila real. \\
\texttt{(cross-cell :sheet var :col col :row fila)} & Referencia a una celda
  concreta de otra hoja, identificada por su variable de hoja, columna y fila.
  Se usa cuando es necesario leer un campo de cada hoja en un bucle de
  generación. \\
\bottomrule
\end{tabular}
\end{center}

\end{document}
